% GENERATED FILE. Edit canonical lessons, then run npm run generate:books.
% canonical-source-hash: fnv1a64:d42ec3df05351f6d
% canonical-lessons: UR-C323-rivayat, UR-C323-rivaj, UR-C323-rasm, UR-C323-manna, UR-C323-nibhana, UR-R323-review-316-323

\chapter{Old Customs}
\label{ch:ur-old-customs}
\hlchaptermodality{323}

\begin{chapteropening}
\textbf{By the end of this chapter:} \emph{I can explain an old custom and how it is kept.}

Complete the last lesson of chapter 323: I can explain an old custom and how it is kept.
\end{chapteropening}

\section[\texorpdfstring{\ur{روایت} (rivāyat)}{روایت (rivāyat)}]{\ur{روایت} (rivāyat) — a tradition}

\label{lesson:UR-C323-rivayat}



\begin{quote}
Before the new one: say the Urdu for a prayer of asking, then the Urdu for charity.
\end{quote}

\subsection*{You'll want to know: \ur{روایت}}
\textbf{\ur{روایت}} — \emph{rivāyat} — \textquotedblleft{}a tradition\textquotedblright{}.

\begin{quote}
\ur{یہ} \ur{ہمارے} \ur{گاؤں} \ur{کی} \ur{پرانی} \ur{روایت} \ur{ہے۔}
\end{quote}

\begin{quote}
\emph{yih hamāre gāoṅ kī purānī rivāyat hai.}
\end{quote}

\begin{quote}
This is an old tradition of our village.
\end{quote}

Feminine.

\begin{cousinweb}
\textbf{\ur{روایت}} (\emph{rivāyat}) is usually traced, through Persian, to Arabic \emph{riwāya}, \textquotedblleft{}a handing on of a report\textquotedblright{}: what is passed down.
\end{cousinweb}

\begin{grammarlens}[title={What you've built}]
One more word for telling what happened.
\end{grammarlens}

\subsection*{Your turn}
\begin{itemize}
\raggedright
  \item \emph{Say it:} \emph{rivāyat}
  \item \emph{Say it:} \emph{rivāyat}, once more
  \item \emph{Say it:} \emph{yih hamāre gāoṅ kī purānī rivāyat hai.} — This is an old tradition of our village.
  \item \emph{Recall:} say the Urdu for a prayer of asking, then the Urdu for charity, then say \emph{rivāyat} again
\end{itemize}

\begin{tcolorbox}[breakable,colback=teal!4,colframe=teal!35!black,title={Before you move on}]
What does \ur{روایت} mean? (\textquotedblleft{}A tradition\textquotedblright{}.) Say it once more.
\end{tcolorbox}
\section[\texorpdfstring{\ur{رواج} (rivāj)}{رواج (rivāj)}]{\ur{رواج} (rivāj) — a custom}

\label{lesson:UR-C323-rivaj}



\begin{quote}
Before the new one: say the Urdu for charity, then the Urdu for a tradition.
\end{quote}

\subsection*{You'll want to know: \ur{رواج}}
\textbf{\ur{رواج}} — \emph{rivāj} — \textquotedblleft{}a custom\textquotedblright{}.

\begin{quote}
\ur{پہلے} \ur{یہ} \ur{رواج} \ur{نہیں} \ur{تھا۔}
\end{quote}

\begin{quote}
\emph{pehle yih rivāj nahīṅ thā.}
\end{quote}

\begin{quote}
This custom did not exist before.
\end{quote}

Masculine.

\begin{cousinweb}
is usually traced, through Persian, to Arabic \emph{rawāj}, \textquotedblleft{}currency, being in use\textquotedblright{}: what people do now. \textbf{\ur{روایت}} is old; a \textbf{\ur{رواج}} may be new.
\end{cousinweb}

\begin{grammarlens}[title={What you've built}]
One more word for telling what happened.
\end{grammarlens}

\subsection*{Your turn}
\begin{itemize}
\raggedright
  \item \emph{Say it:} \emph{rivāj}
  \item \emph{Say it:} \emph{rivāj}, once more
  \item \emph{Say it:} \emph{pehle yih rivāj nahīṅ thā.} — This custom did not exist before.
  \item \emph{Recall:} say the Urdu for charity, then the Urdu for a tradition, then say \emph{rivāj} again
\end{itemize}

\begin{tcolorbox}[breakable,colback=teal!4,colframe=teal!35!black,title={Before you move on}]
What does \ur{رواج} mean? (\textquotedblleft{}A custom\textquotedblright{}.) Say it once more.
\end{tcolorbox}
\section[\texorpdfstring{\ur{رسم} (rasm)}{رسم (rasm)}]{\ur{رسم} (rasm) — a rite; a ceremony}

\label{lesson:UR-C323-rasm}



\begin{quote}
Before the new one: say the Urdu for a tradition, then the Urdu for a custom.
\end{quote}

\subsection*{You'll want to know: \ur{رسم}}
\textbf{\ur{رسم}} — \emph{rasm} — \textquotedblleft{}a rite; a ceremony\textquotedblright{}.

\begin{quote}
\ur{شادی} \ur{سے} \ur{پہلے} \ur{مہندی} \ur{کی} \ur{رسم} \ur{ہوئی۔}
\end{quote}

\begin{quote}
\emph{shādī se pehle mehndī kī rasm huī.}
\end{quote}

\begin{quote}
Before the wedding there was the henna ceremony.
\end{quote}

Feminine.

\begin{cousinweb}
\textbf{\ur{رسم}} (\emph{rasm}) is usually traced to Persian \emph{rasm}, \textquotedblleft{}a usage, a custom\textquotedblright{}, a word Persian took from Arabic. It is feminine: \textbf{\ur{رسم} \ur{ہوئی}}.
\end{cousinweb}

\begin{grammarlens}[title={What you've built}]
One more word for telling what happened.
\end{grammarlens}

\subsection*{Your turn}
\begin{itemize}
\raggedright
  \item \emph{Say it:} \emph{rasm}
  \item \emph{Say it:} \emph{rasm}, once more
  \item \emph{Say it:} \emph{shādī se pehle mehndī kī rasm huī.} — Before the wedding there was the henna ceremony.
  \item \emph{Recall:} say the Urdu for a tradition, then the Urdu for a custom, then say \emph{rasm} again
\end{itemize}

\begin{tcolorbox}[breakable,colback=teal!4,colframe=teal!35!black,title={Before you move on}]
What does \ur{رسم} mean? (\textquotedblleft{}A rite; a ceremony\textquotedblright{}.) Say it once more.
\end{tcolorbox}
\section[\texorpdfstring{\ur{ماننا} (mānnā)}{ماننا (mānnā)}]{\ur{ماننا} (mānnā) — to accept; to agree to}

\label{lesson:UR-C323-manna}



\begin{quote}
Before the new one: say the Urdu for a custom, then the Urdu for a rite.
\end{quote}

\subsection*{You'll want to know: \ur{ماننا}}
\textbf{\ur{ماننا}} — \emph{mānnā} — \textquotedblleft{}to accept; to agree to\textquotedblright{}.

\begin{quote}
\ur{آخر} \ur{اس} \ur{نے} \ur{میری} \ur{بات} \ur{مانی۔}
\end{quote}

\begin{quote}
\emph{ākhir us ne merī bāt mānī.}
\end{quote}

\begin{quote}
In the end he accepted what I said.
\end{quote}

Past, with \ur{نے} (he or she …): \ur{اس} \ur{نے} \ur{سچ} \ur{مانا}, \emph{us ne sach mānā}.

\begin{cousinweb}
\textbf{\ur{ماننا}} (\emph{mānnā}) is inherited: CDIAL lists the Hindi cousin \emph{mānnā}, \textquotedblleft{}to respect, to obey\textquotedblright{}, under Sanskrit \emph{mányatē}, \textquotedblleft{}thinks\textquotedblright{} (CDIAL 9857). \textbf{\ur{بات} \ur{ماننا}} is to accept what someone says.
\end{cousinweb}

\begin{grammarlens}[title={What you've built}]
One more word for telling what happened.
\end{grammarlens}

\subsection*{Your turn}
\begin{itemize}
\raggedright
  \item \emph{Say it:} \emph{mānnā}
  \item \emph{Say it:} \emph{mānnā}, once more
  \item \emph{Say it:} \emph{ākhir us ne merī bāt mānī.} — In the end he accepted what I said.
  \item \emph{Recall:} say the Urdu for a custom, then the Urdu for a rite, then say \emph{mānnā} again
\end{itemize}

\begin{tcolorbox}[breakable,colback=teal!4,colframe=teal!35!black,title={Before you move on}]
What does \ur{ماننا} mean? (\textquotedblleft{}To accept; to agree to\textquotedblright{}.) Say it once more.
\end{tcolorbox}
\section[\texorpdfstring{\ur{نبھانا} (nibhānā)}{نبھانا (nibhānā)}]{\ur{نبھانا} (nibhānā) — to keep up; to carry through}

\label{lesson:UR-C323-nibhana}



\begin{quote}
Before the new one: say the Urdu for a rite, then the Urdu for to accept.
\end{quote}

\subsection*{You'll want to know: \ur{نبھانا}}
\textbf{\ur{نبھانا}} — \emph{nibhānā} — \textquotedblleft{}to keep up; to carry through\textquotedblright{}.

\begin{quote}
\ur{ہم} \ur{نے} \ur{یہ} \ur{روایت} \ur{نبھائی۔}
\end{quote}

\begin{quote}
\emph{ham ne yih rivāyat nibhāī.}
\end{quote}

\begin{quote}
We kept up this tradition.
\end{quote}

Past, with \ur{نے} (he or she …): \ur{اس} \ur{نے} \ur{رسم} \ur{نبھائی}, \emph{us ne rasm nibhāī}.

\begin{cousinweb}
\textbf{\ur{نبھانا}} (\emph{nibhānā}) is the causing form of \emph{nibhnā}, \textquotedblleft{}to last, to be carried through\textquotedblright{}, which CDIAL lists for Hindi under Sanskrit \emph{nírvahati} (CDIAL 7397): to make a custom or a promise last.
\end{cousinweb}

\begin{grammarlens}[title={What you've built}]
One more word for telling what happened.
\end{grammarlens}

\subsection*{Your turn}
\begin{itemize}
\raggedright
  \item \emph{Say it:} \emph{nibhānā}
  \item \emph{Say it:} \emph{nibhānā}, once more
  \item \emph{Say it:} \emph{ham ne yih rivāyat nibhāī.} — We kept up this tradition.
  \item \emph{Recall:} say the Urdu for a rite, then the Urdu for to accept, then say \emph{nibhānā} again
\end{itemize}

\begin{tcolorbox}[breakable,colback=teal!4,colframe=teal!35!black,title={Before you move on}]
What does \ur{نبھانا} mean? (\textquotedblleft{}To keep up; to carry through\textquotedblright{}.) Say it once more.
\end{tcolorbox}
\section[(dialogue)]{Looking back: from Escape from the Fort to Old Customs}

\label{lesson:UR-R323-review-316-323}



\begin{quote}
Say the Urdu for a prisoner and a guard kept.
\end{quote}

\subsection*{Your turn}
Say these aloud:

\begin{itemize}
\raggedright
  \item a prisoner, a guard kept, a tunnel, to slip away, to crawl
  \item the festival of Eid, the new crescent moon, henna, to glitter, festive bustle
  \item the set prayer, money given at Eid, sweet vermicelli, congratulations, to cling
  \item the spring kite festival, a kite string, a spool, to fly a kite, to be cut
  \item fireworks, a firecracker, to sprinkle, to drench, to sway
  \item to swing, a procession, a crowd, gathered together, a spectacle
  \item a fast, the meal before dawn, the food that breaks a fast, a prayer of asking, charity
  \item a tradition, a custom, a rite, to accept, to keep up
\end{itemize}

\begin{tcolorbox}[breakable,colback=teal!4,colframe=teal!35!black,title={Before you move on}]
A prisoner? (\textbf{\ur{قیدی}}, \emph{qaidī}.) A guard kept? (\textbf{\ur{پہرہ}}, \emph{pahrā}.) A tunnel? (\textbf{\ur{سرنگ}}, \emph{surang}.) To slip away? (\textbf{\ur{کھسکنا}}, \emph{khisaknā}.) To crawl? (\textbf{\ur{رینگنا}}, \emph{rengnā}.) The festival of Eid? (\textbf{\ur{عید}}, \emph{īd}.) The new crescent moon? (\textbf{\ur{ہلال}}, \emph{hilāl}.) Henna? (\textbf{\ur{مہندی}}, \emph{mehndī}.) To glitter? (\textbf{\ur{جگمگانا}}, \emph{jagmagānā}.) Festive bustle? (\textbf{\ur{رونق}}, \emph{raunaq}.) The set prayer? (\textbf{\ur{نماز}}, \emph{namāz}.) Money given at Eid? (\textbf{\ur{عیدی}}, \emph{īdī}.) Sweet vermicelli? (\textbf{\ur{سویاں}}, \emph{sivayāṅ}.) Congratulations? (\textbf{\ur{مبارک} \ur{باد}}, \emph{mubārak bād}.) To cling? (\textbf{\ur{لپٹنا}}, \emph{lipaṭnā}.) The spring kite festival? (\textbf{\ur{بسنت}}, \emph{basant}.) A kite string? (\textbf{\ur{ڈور}}, \emph{ḍor}.) A spool? (\textbf{\ur{چرخی}}, \emph{charkhī}.) To fly a kite? (\textbf{\ur{اڑانا}}, \emph{uṛānā}.) To be cut? (\textbf{\ur{کٹنا}}, \emph{kaṭnā}.) Fireworks? (\textbf{\ur{آتش} \ur{بازی}}, \emph{ātish bāzī}.) A firecracker? (\textbf{\ur{پٹاخا}}, \emph{paṭākhā}.) To sprinkle? (\textbf{\ur{چھڑکنا}}, \emph{chhiṛaknā}.) To drench? (\textbf{\ur{بھگونا}}, \emph{bhigonā}.) To sway? (\textbf{\ur{جھومنا}}, \emph{jhūmnā}.) To swing? (\textbf{\ur{جھولنا}}, \emph{jhūlnā}.) A procession? (\textbf{\ur{جلوس}}, \emph{julūs}.) A crowd? (\textbf{\ur{ہجوم}}, \emph{hujūm}.) Gathered together? (\textbf{\ur{اکٹھا}}, \emph{ikaṭṭhā}.) A spectacle? (\textbf{\ur{تماشا}}, \emph{tamāshā}.) A fast? (\textbf{\ur{روزہ}}, \emph{rozā}.) The meal before dawn? (\textbf{\ur{سحری}}, \emph{sahrī}.) The food that breaks a fast? (\textbf{\ur{افطاری}}, \emph{iftārī}.) A prayer of asking? (\textbf{\ur{دعا}}, \emph{duā}.) Charity? (\textbf{\ur{خیرات}}, \emph{khairāt}.) A tradition? (\textbf{\ur{روایت}}, \emph{rivāyat}.) A custom? (\textbf{\ur{رواج}}, \emph{rivāj}.) A rite? (\textbf{\ur{رسم}}, \emph{rasm}.) To accept? (\textbf{\ur{ماننا}}, \emph{mānnā}.) To keep up? (\textbf{\ur{نبھانا}}, \emph{nibhānā}.)
\end{tcolorbox}
